\chapter{Preface}

This is not Manuel Bronstein's book. It is an attempt to write down, in one
place and in one notation, the material that the second volume of his
\emph{Symbolic Integration} was to have covered: the integration of algebraic
functions, and of functions built from algebraic and transcendental ones
together.

\section*{What is known of the second volume}

That a second volume was planned is on record in the first. Its preface says:
\begin{quote}
The reader will notice that several topics in symbolic integration are missing
from this book, the main one being the integration of algebraic functions.
Including algorithms for integrating algebraic and mixed
algebraic-transcendental functions would however easily double the size of
this book, as well as increase the mathematical prerequisites, since those
algorithms require prior familiarity with algebraic curves and functions. We
have thus decided to cover algebraic functions in a second volume, which will
hopefully appear in the near future. \cite[Preface to the First
Edition]{bronstein2005}
\end{quote}
Bronstein died on 6 June 2005. S.~A. Abramov's obituary records what became of
the plan: ``The first volume was written and went through two editions at the
Springer publishing house in 1997 and 2004. The second volume remained
uncompleted''~\cite{abramov2006}.

That is all we know. No manuscript, draft, outline or table of contents of the
second volume is known to us, and nothing in this book is taken from one. The
choice of topics, their order and their treatment are a reconstruction, made
from what the first volume says the second was for, from the literature it
sends its reader to, and from Bronstein's own shorter writings on the subject.
Where this book and the book he would have written differ, which must be
almost everywhere, the fault is not his.

\section*{What this book is}

It is a companion to the first volume, meant to be read after it. It keeps
that book's notation and the shape of its chapters --- the theory, then the
algorithm as pseudocode, then an example worked through it, then exercises ---
and it cites the first volume by section for whatever is proved there.

It is assembled from the literature and contains no new research. Where the
sources give a complete algorithm, the book gives it, with a proof or with a
citation to the place where the proof is. Where they give a sketch, the book
says that it is a sketch. Where a problem is open, the book states it as open
and does not fill the gap with an algorithm of its own.
Appendix~\ref{app:standing} lists every result with its standing.

Some of the material is later than 2005, and could not have been in any book
of Bronstein's. It is included because a reader today needs it, and it is
marked with a dagger\later.
